\section{本章小结}
\label{sec:simplex-summary}

本章开头的问题是怎样在有限但极多的顶点之间找到最优的一个。单纯形法的回答是：用基来表示顶点，用约化成本判断哪条边值得走，用比值检验决定走多远，用对偶候选写出最优性证书。工厂的例子经过三次换基，从$(0,0)$走到$(4,2)$，并得到资源的对偶价格$(2,1,0)$。

算法的每个环节都有明确的条件。起步需要可行基，没有时用人工变量，两阶段法或符号化的大M法把它们排除；有限的$M$要超过对偶价格，这一点事先并不知道。退化基使步长为零，同一点对应多个基，循环可能发生，Bland规则或字典序规则在精确算术下保证有限终止。实现上修正单纯形法只维护基矩阵的因子，浮点运算要求小心选主元、定期重新分解，并在结束时重新核查。

算法的结果总是带着证据，表~\ref{tab:simplex-outcomes}汇总了三种停止结果。

\begin{BookInlineTable}
 \centering
 \begin{tabular}{lll}
  \toprule
  结果 & 检验 & 证据\\
  \midrule
  最优 & 可行基上$\bm{s}\geq\bm{0}$ & 原始解和对偶解，目标相等\\
  无界 & $s_q<0$且$\bm{u}\leq\bm{0}$ & 可行点和方向$\bm{d}$，$\bm{c}^\top\bm{d}<0$\\
  不可行 & 第一阶段$\phi^\star>0$，& $\bm{y}$，满足$\bm{A}^\top\bm{y}\leq\bm{0}$，$\bm{b}^\top\bm{y}>0$\\
  & 或对偶单纯形法无候选 & （后者取$\bm{y}=-\bm{\rho}$）\\
  \bottomrule
 \end{tabular}
 \caption{单纯形法的三种停止结果及其证据。}
 \label{tab:simplex-outcomes}
\end{BookInlineTable}

得到最优基以后，还能继续使用它。基的可行性和最优性分别管着右端和目标系数，固定基的参数分析读出每个参数的有效范围和影子价格；越出范围后，右端的变化交给对偶单纯形法，目标系数的变化交给原始单纯形法，只需少量换基。这种复用是单纯形法在整数规划等反复求解场合的重要优势。

有限终止不等于快。Klee--Minty构造说明Dantzig规则的换基次数最坏可以指数增长，实践中换基次数虽然通常不大，却没有无条件的理论保证。第~\ref{chap:lp-interior-point}章将研究另一条路线：不沿边界走，而从可行域内部逼近最优解。
